\documentclass[11pt,a4paper]{article}
\usepackage[T1]{fontenc}
\usepackage[utf8]{inputenc}
\usepackage{lmodern,amsmath,amssymb}
\usepackage[margin=25mm]{geometry}
\usepackage[hidelinks]{hyperref}
\hypersetup{pdftitle={What would unblock the confinement-scale region, and what was checked and fails},pdfauthor={navstokgap research working notes}}
\newcommand{\tightlist}{\setlength{\itemsep}{0pt}\setlength{\parskip}{0pt}}
\setlength{\emergencystretch}{3em}
\setcounter{secnumdepth}{0}
\begin{document}
\hypertarget{what-would-unblock-the-confinement-scale-region-and-what-was-checked-and-fails}{%
\section{What would unblock the confinement-scale region, and what was
checked and
fails}\label{what-would-unblock-the-confinement-scale-region-and-what-was-checked-and-fails}}

The map of \href{mass-gap-position.md}{the position note} leaves one
region untouched by every tool with explicit constants: the confinement
scale, where the effective coupling is of order one and the effective
theory must be shown to be mixing uniformly in the volume. This note
records the outcome of a review of every rigorous technique known to the
author that might enter that region by written derivation, so that later
sessions do not re-derive the same dead ends. The conclusion is that
exactly two things would unblock it, and neither is available here:

\begin{enumerate}
\def\labelenumi{\arabic{enumi}.}
\tightlist
\item
  \textbf{A correlation inequality for \(SU(3)\) lattice gauge theory}
  of Griffiths--Ginibre type, giving monotonicity in the coupling of the
  quantities that control mixing. This is what makes the Ising model's
  intermediate regime accessible and what the abelian gauge theories
  have; for \(SU(N)\) none is known, and Section 1 says where the
  standard constructions break.
\item
  \textbf{A computer-assisted verification of the Dobrushin--Shlosman
  block criterion} for the Wilson interaction, or for the effective
  interaction at the confinement scale, as specified in
  \href{dobrushin-uniqueness-wilson.md}{the Dobrushin note} §4. This is
  numerics of a size beyond present practice and outside this
  repository's rules.
\end{enumerate}

Everything else checked reduces to one of these or fails for a stated
reason. References at metadata level; nothing promoted.

\hypertarget{correlation-inequalities-where-they-break-for-sun}{%
\subsection{\texorpdfstring{1. Correlation inequalities: where they
break for
\(SU(N)\)}{1. Correlation inequalities: where they break for SU(N)}}\label{correlation-inequalities-where-they-break-for-sun}}

\emph{Ising and abelian gauge theories.} Griffiths' inequalities in
Ginibre's general form (Commun. Math. Phys. 16 (1970) 310) need a cone
of functions on the single-site space, closed under products, such that
the duplicate-variables integral
\(\int\!\!\int\prod_i\big(f_i(x)-f_i(y)\big)\,d\nu(x)\,d\nu(y)\ge0\)
holds. For \(\pm1\) spins and for \(U(1)\) with the cosines this holds,
so the Wilson loops of \(U(1)\) lattice gauge theory are monotone in
\(\beta_W\) and the phase structure can be squeezed from both sides; the
sharpness theorem of Aizenman, Barsky and Fernández (J. Stat. Phys. 47
(1987) 343) is the deepest form of this control for Ising-type models.
The Coulomb phase of \(U(1)\) (Guth; Fröhlich--Spencer) uses the dual
representation, whose weights are positive.

\emph{Non-abelian.} Three obstructions, each sufficient on its own.

\begin{itemize}
\tightlist
\item
  The dual representation has signs. Integrating the links of a
  character expansion produces recoupling coefficients, which are not
  nonnegative for \(SU(2)\) and \(SU(3)\), so the surface gas is not a
  positive measure and no monotonicity in \(\beta_W\) follows from it.
  The plaquette coefficients themselves are nonnegative
  (\href{wilson-strong-coupling-explicit.md}{Wilson note} §1), which is
  why the strong-coupling expansion is a positive gas; the signs appear
  at the vertices where more than two surfaces meet.
\item
  The \(O(N)\) analogy fails already at \(N=4\). \(SU(2)\simeq S^3\),
  and \(\tfrac12\operatorname{tr}(UV^\dagger)\) is the Euclidean inner
  product of unit quaternions, so the plaquette term is a degree-four
  polynomial in four \(O(4)\) spins; Ginibre's construction covers
  \(O(2)\) and specific \(O(N)\) interactions but not this one, and
  \(SU(3)\) has no such real form at all.
\item
  The interaction is not ferromagnetic in any sign sense: for a complex
  defining representation \(\operatorname{Re}\operatorname{tr}U_p\) is
  not a product of functions on the links with definite sign properties.
\end{itemize}

\hypertarget{reflection-positivity-tools}{%
\subsection{2. Reflection-positivity
tools}\label{reflection-positivity-tools}}

\emph{Chessboard estimates} (Fröhlich, Israel, Lieb and Simon, Commun.
Math. Phys. 62 (1978) 1) hold for the Wilson measure and bound the
probability of a local event by a free-energy difference per block,
uniformly in the volume and at every coupling. They give rigorous
large-deviation bounds for bad blocks without an expansion. They do not
give mixing on the good blocks, so they are an ingredient of a
verification and not a substitute for it.

\emph{Infrared bounds} (Fröhlich, Simon and Spencer, Commun. Math. Phys.
50 (1976) 79) need Gaussian domination, \(Z(h)\le Z(0)\) for shifts
\(h\) of the field, which uses the translation structure of a flat
target space or of \(U(1)\); on a non-abelian group manifold the shift
is not a symmetry and the domination fails. This is consistent with
confinement: an infrared bound would prove a massless phase.

\emph{Reflection positivity proves transitions} (Peierls-type arguments
via chessboard) and, through the transfer matrix, converts decay into a
gap (\href{dobrushin-uniqueness-wilson.md}{Dobrushin note} §3). It does
not prove the absence of a transition.

\hypertarget{hamiltonian-finite-size-criteria}{%
\subsection{3. Hamiltonian finite-size
criteria}\label{hamiltonian-finite-size-criteria}}

Knabe's criterion (J. Stat. Phys. 52 (1988) 627) and its descendants
bound the gap of an infinite frustration-free Hamiltonian by the gap on
finite blocks. The Kogut--Susskind Hamiltonian is not frustration-free:
its ground state is not a ground state of the electric and magnetic
terms separately, and no equivalent frustration-free form with the same
gap is known. Ground-state clustering also does not imply a gap for a
general Hamiltonian, while Euclidean decay in time does, through the
positive transfer matrix; this is why the decay-to-gap step in this
programme is always taken on the Euclidean side.

\hypertarget{exact-transformations}{%
\subsection{4. Exact transformations}\label{exact-transformations}}

\emph{Flow conjugation} is exact and preserves the spectrum at the same
coupling (\href{flow-conjugation-truncation.md}{flow-conjugation note}).
\emph{Electric--magnetic duality} is exact for abelian groups and maps
strong to weak coupling; for \(SU(N)\) the dual is the signed surface
gas of Section 1. \emph{Gauge fixing} changes the specification, so the
Dobrushin coefficient can be reduced by removing fixed links from the
neighbour count, by at most the fraction of links on a maximal tree,
about one quarter in four dimensions; this moves \(444\) to about
\(330\) and touches nothing else.

\hypertarget{other-constructions}{%
\subsection{5. Other constructions}\label{other-constructions}}

\emph{Large \(N\)} has \(1/N\) as a small parameter; the goal is
\(SU(3)\). \emph{Stochastic quantization} (the Langevin dynamics of the
gauge field) constructs the two- and three-dimensional theories locally
in time and says nothing about the gap; four dimensions is not
accessible. \emph{Small volume} relocates the crossover from the
coupling to the volume variable: the small-volume theory has an explicit
gap \(\delta_1g^{2/3}\hbar c/L\), and the passage to \(L\Lambda\gg1\) is
the same confinement-scale problem.

\hypertarget{the-confinement-scale-as-a-single-lattice-theory}{%
\subsection{6. The confinement scale as a single lattice
theory}\label{the-confinement-scale-as-a-single-lattice-theory}}

After the small-field renormalization has run to the scale where the
effective coupling is of order one, the problem is one lattice gauge
theory at that scale with a bounded but unknown effective interaction,
to be shown mixing uniformly in the volume. Beyond that scale the block
variables decorrelate and the effective interaction approaches the
trivial fixed point, inside the completely analytical set. The whole
difficulty is therefore concentrated at one scale and one effective
interaction, which is the sharpest form of the statement that the
problem has no small parameter, and which is exactly what the block
verification of item 2 addresses.

\hypertarget{consequence-for-state}{%
\subsection{7. Consequence for STATE}\label{consequence-for-state}}

The review is closed. By written derivation, the confinement-scale
region is entered by nothing known; the two unblocking routes are a
correlation inequality for \(SU(3)\), which would be new mathematics,
and the computer-assisted block verification, which is a decision about
the repository's rules and a large project. Constant-tightening on
either side of the region remains available and does not touch it.
\end{document}
